% from ega3-1-fr.tex
\oldpage[0\textsubscript{III}]{48}
Prenons alors
\[
Q_{i+1}=Q'_{i+1}\times Q''_{i+1},\qquad
d_{i+1}=d'_{i+1}+d''_{i+1},\qquad
u_{i+1}=u'_{i+1}+u''_{i+1}.
\]
Comme $d_{i+1}(Q_{i+1})=d'_{i+1}(Q'_{i+1})=N_i\subset Z_i(Q_\bullet)$, on a
$d_i\circ d_{i+1}=0$ et, avec les notations usuelles,
$B_i(Q_\bullet)=N_i$, ce qui vérifie $(\mathrm I_{i+1})$. La commutativité
des diagrammes (\hyperref[0.11.9.1.3-fr]{11.9.1.3}) et
(\hyperref[0.11.9.1.4-fr]{11.9.1.4}) montre que
$d_{i+1}\circ u_{i+1}=u_i\circ d_{i+1}$. Par définition de $N_i$, le
morphisme
\[
H_i(Q_\bullet)=Z_i(Q_\bullet)/N_i\longrightarrow H_i(P_\bullet)
\]
déduit du système des $u_k$ ($k\leq i+1$) est le morphisme déduit de $v_i$
par passage aux quotients, donc c'est un isomorphisme puisque $v_i$ est un
épimorphisme, d'où $(\mathrm{II}_{i+1})$. Enfin, on a
$Q''_{i+1}\subset Z_{i+1}(Q_\bullet)$ par définition; le choix de
$u''_{i+1}$ montre que le morphisme
\[
v_{i+1}:Z_{i+1}(Q_\bullet)\longrightarrow Z_{i+1}(P_\bullet)
\longrightarrow H_{i+1}(P_\bullet)
\]
est un épimorphisme, sa restriction à $Q''_{i+1}$ l'étant déjà, d'où
$(\mathrm{III}_{i+1})$. La récurrence peut donc se poursuivre, et la
proposition est démontrée.
\end{proof}

\begin{corollary}[11.9.2]
\label{0.11.9.2-fr}
Soient $A$ un anneau noethérien (non nécessairement commutatif),
$P_\bullet=(P_i)_{i\in\mathbb Z}$ un complexe de $A$-modules à droite. On
suppose que les $H_i(P_\bullet)$ sont des $A$-modules de type fini et qu'il
existe $d$ tel que $H_i(P_\bullet)=0$ pour $i<d$. Alors il existe un complexe
$Q_\bullet=(Q_i)_{i\in\mathbb Z}$ formé de $A$-modules à droite libres de
rang fini, tel que $Q_i=0$ pour $i<d$, et un homomorphisme
$u:Q_\bullet\to P_\bullet$ de complexes, tel que l'homomorphisme
$H_\bullet(Q_\bullet)\to H_\bullet(P_\bullet)$ correspondant à $u$ soit
bijectif.
\end{corollary}

\begin{proof}
On applique (\hyperref[0.11.9.1-fr]{11.9.1}) en prenant pour $\mathcal C$ la
catégorie des $A$-modules à droite, pour $K'$ (resp. $K''$) l'ensemble des
$A$-modules de type fini (resp. l'ensemble des $A$-modules libres de rang
fini); la vérification des conditions (i), (ii) et (iii) de
(\hyperref[0.11.9.1-fr]{11.9.1}) est immédiate, compte tenu de l'hypothèse
que $A$ est noethérien.
\end{proof}

\begin{env}[11.9.3]
\label{0.11.9.3-fr}
\emph{Remarques}.
\begin{enumerate}
\item[{\rm(i)}] Sous les conditions de
(\hyperref[0.11.9.2-fr]{11.9.2}), supposons en outre que les $P_i$ soient des
$A$-modules à droite plats. Alors, pour tout $A$-module à gauche $M$,
l'homomorphisme de complexes
\[
u\otimes 1:Q_\bullet\otimes_A M\longrightarrow P_\bullet\otimes_A M
\]
définit encore un isomorphisme
$H_\bullet(Q_\bullet\otimes_A M)\simeq
H_\bullet(P_\bullet\otimes_A M)$ de l'homologie comme nous le verrons au
chap.~III.
\item[{\rm(ii)}] La conclusion de (\hyperref[0.11.9.2-fr]{11.9.2}) n'est
plus nécessairement exacte lorsqu'on ne suppose pas $A$ noethérien; en
effet, en l'appliquant à un complexe réduit à $0$ sauf pour un seul terme, on
en conclurait que tout $A$-module à gauche de type fini admet une résolution
par des modules libres de type fini, ce qui n'est pas vrai en général
(cf. Bourbaki, \emph{Alg. comm.}, chap.~I, \S~2, exerc.~6).

Toutefois, au lieu de supposer $A$ noethérien, on peut supposer seulement que
les $H_i(P_\bullet)$ ont une $\infty$-présentation finie (cf. chap.~IV).
\end{enumerate}
\end{env}

\subsection{Caractéristique d'Euler-Poincaré d'un complexe de modules de
longueur finie}
\label{subsection:0.11.10-fr}

\begin{env}[11.10.1]
\label{0.11.10.1-fr}
Soient $A$ un anneau (non nécessairement commutatif),
\[
\label{0.11.10.1.1-fr}
\tag{11.10.1.1}
M^\bullet:0\longrightarrow M^0\longrightarrow M^1
\longrightarrow\cdots\longrightarrow M^n\longrightarrow0
\]
un complexe de $A$-modules à gauche de longueur finie. On appelle
\emph{caractéristique d'Euler-Poincaré} de ce complexe le nombre
